\documentclass{article}
\usepackage{natbib,graphics,amsmath,kmq2,ifthen,amssymb,array,
  indentfirst}
\usepackage{rotating}
 
\renewcommand{\a}{\boldsymbol{\alpha}}
\renewcommand{\b}{\boldsymbol{\beta}}
\newcommand{\e}{\boldsymbol{\eta}}
\newcommand{\g}{\boldsymbol{\gamma}}
\newcommand{\m}{\boldsymbol{\mu}}
\newcommand{\mm}{\mathbf{m}}
\newcommand{\p}{\boldsymbol{\psi}}
\newcommand{\ph}{\boldsymbol{\phi}}
\renewcommand{\S}{\boldsymbol{\Sigma}}
\newcommand{\T}{\boldsymbol{\Theta}}
\renewcommand{\t}{\boldsymbol{\theta}}
\renewcommand{\L}{\boldsymbol{\Lambda}}
\renewcommand{\l}{\boldsymbol{\lambda}}
\newcommand{\x}{\mathbf{x}}
\newcommand{\X}{\mathbf{X}}
\newcommand{\y}{\mathbf{y}}
\newcommand{\w}{\mathbf{w}}
\newcommand{\W}{\mathbf{W}}
\renewcommand{\u}{\mathbf{u}}
\newcommand{\U}{\mathbf{U}}
\newcommand{\z}{\mathbf{z}}
\newcommand{\Z}{\mathbf{Z}}
\newcommand{\V}{\mathbf{V}}
\newcommand{\E}{\mathbf{E}}
\newcommand{\Y}{\mathbf{Y}}
\newcommand{\F}{\mathbf{F}}
\renewcommand{\I}{\mathbb{I}}

\newtheorem{property}{Property}
%\newtheorem{assumption}{Assumption}
\newtheorem{assumption}{Effect Maximizing Assumption}




%\newcommand{\ob}[1]{\textcolor[rgb]{0.65, 0.65, 0.65}{{#1}}}
\newcommand{\ob}[1]{\textcolor[rgb]{1, 0, 0}{{#1}}}


%% === hyperref options ===
\usepackage{color}
\usepackage[pdftex, bookmarksopen=true, 
bookmarksnumbered=true, linkcolor=webred]{hyperref}
\definecolor{webgreen}{rgb}{0, 0.5, 0}
\definecolor{webblue}{rgb}{0, 0, 0.5}
\definecolor{webred}{rgb}{0.5, 0, 0}


\newcommand{\indep}{{\bot\negthickspace\negthickspace\bot}}
\newcommand{\notindep}{{\bot\negthickspace\negthickspace\bot}{\negthinspace
  \negmedspace \negthickspace \negthickspace \diagup
  \;}}
\newcommand{\notone}{1 \negthinspace \negthickspace \negthickspace \diagup} 
\newcommand{\notnegone}{-1 \negthickspace \negthickspace \negthickspace \diagup}
\newcommand{\notzero}{0 \negthinspace \negthickspace \negthickspace \diagup}


\title{Gov 1000/2000/2000e Causality Lecture Notes}

\author{Adam Glynn}
\date{}

\begin{document}
%\maketitle

\double

Suppose we allow an interaction between the treatment and the conditioning variable
\begin{align*}
E[Y|T=t,X=x] = \beta_0 + \beta_1 \cdot t + \beta_2 \cdot x+ \beta_3 \cdot t \cdot x
\end{align*}
In this model, $\beta_1$ can \emph{no longer} be interpreted predictively as the difference between two conditional expectations. Instead, we predictively interpret $(\beta_1 + \beta_3 \cdot x)$ as the difference between two conditional expectations $E[Y|T=t+1,X=x]-E[Y|T=t,X=x]$ which depends on the value of $X=x$. This can be most easily seen by rearranging terms in the conditional expectation:
\begin{align*}
E[Y|T=t,X=x] = \beta_0 + \bs{(\beta_1 + \beta_3 \cdot x)} \cdot t + \beta_2 \cdot x
\end{align*}
Conditional exchangeability allows us to interpret $(\beta_1 + \beta_3 \cdot x)$ causally (in terms of potential outcomes): 
\begin{align}
(\beta_1 + \beta_3 \cdot x) &= E[Y|T=t+1,X=x]-E[Y|T=t,X=x] \text{ for all } t,x \notag \\
&= E[Y|T=t+1,X=x]-E[Y|T=t,X=x] \text{ for all } t \notag \\
&= E[y^{t+1}|T=t+1,X=x]-E[y^{t}|T=t,X=x] \text{ for all } t \notag\\
&= E[y^{t+1}|X=x]-E[y^{t}|X=x] \text{ for all } t \notag \\
&= E[y^{t+1}-y^{t}|X=x]\ \text{ for all } t \label{eq:condeff}
\end{align}

However, although the average treatment effect conditional on $X=x$ ($E[y^{t+1}-y^{t}|X=x]$) may be of substantive interest,  the law of total expectation, and our linearity assumption allow us to interpret $(\beta_1 + \beta_3 \cdot E[X])$ as the average treatment effect. 
\begin{align*}
\beta_1 + \beta_3 \cdot E_X\left\{X \right\} &= E_X\left\{\beta_1 + \beta_3 \cdot X \right\} \\
&= E_X\left\{E[Y|T=t+1,X=x]-E[Y|T=t,X=x]\right\}\\
&= E_X\left\{E[y^{t+1}|T=t+1,X=x]-E[y^{t}|T=t,X=x]\right\} \\
&= E_X\left\{E[y^{t+1}|X=x]-E[y^{t}|X=x]\right\}\\
&= E_X\left\{E[y^{t+1}-y^{t}|X=x]\right\} \\
&= E[y^{t+1}-y^{t}] 
\end{align*}

The variance is more difficult. We know that conditional on $X=x$, we have a marginally randomized experiment such that the standard regression estimate of variance for $(\hat{\beta}_1 + \hat{\beta}_3 \cdot x)$ will tend to be conservative. However, the variance for the estimated ATE will be a bit more complicated
\begin{align*}


\end{align*}



%For some applications, we will only be interested combinations of the proportions of response types described in the previous section. In the program evaluation setting, we often want to know the average effect of the program on the participants. This average effect, sometimes called the Average Treatment effect on the Treated (ATT). The ATT is just the average of the ITEs for all individuals in the population that received the treatment.In the social pressure example, ATT describes the proportion caused to vote by social pressure minus the proportion caused to not vote by social pressure, and hence this is the increase in voter turnout among those that received the mailing that is due to the mailing.\footnote{With a dichotomous outcome, ATT can always be described as the difference between the proportion caused among the treated minus the proportion prevented among the treated.}  

%If instead of evaluating an existing program, we consider extending a program to the control units, then we might like to know the average effect of the program among the non-participants. This average effect is sometimes called the Average Treatment effect on the Controls (ATC). In the social pressure example, ATC describes the proportion that would have been caused to vote by the mailing minus the proportion that would have been caused to not vote by the mailing, and hence this is the increase in voter turnout among the controls that we would have seen if they had been sent the mailing.\footnote{With a dichotomous outcome, ATC can always be  can be described as the proportion caused minus the proportion prevented among the control units.} 

%Finally, if instead we want to compare the average effect of the policy that treats all units in the population (send the mailing to everyone) to the policy that treats no units in the population (send the mailing to no one), then we need only average ATT and ATC according to the probability of treatment.
%\begin{align*}
%ATE = ATT \cdot Pr(T=1) + ATC \cdot Pr(T=0)
%\end{align*}
%In the social pressure example, this would involve weighting ATT according to the number of individuals that received the mailing and weighting ATC according to the number of individuals that did not receive the mailing. This is known as the Average Treatment Effect (ATE), and can be alternatively described as the average of the ITEs for the entire population ($\frac{1}{N}\sum_{i=1}^N y_i^1- y_i^0$ where $N$ is the population size). See GH pg. 173 for an explanation with a slightly different focus. 

%\subsection{Identification of the Average Treatment Effect with Randomized Experiments}

%One reason for the disciplinary focus on average effects is that they can be estimated with randomized experiments and we only need to assume SUTVA (and even this can be relaxed in some cases). Furthermore, if we understand why randomized experiments allow us to estimate average effects, we can use this knowledge to inform our interpretations of non-randomized studies. 

%\begin{quote}
%{\it A good way to do research is to imagine the experiment that we would like to conduct, and then assess how close our data get us to that experiment.}
%\end{quote}
%\begin{center} I'm paraphrasing someone, but I can't remember who\end{center}

%To understand why randomization allows the estimation of average effects, we must utilize some mathematical notation.  
%Recall that the ATE can be written as, $\frac{1}{N}\sum_{i=1}^N (y_i^1- y_i^0)$. Because this is an average, it can be re-written in the following manner:
%\begin{align*}
%\frac{1}{N}\sum_{i=1}^N (y_i^1- y_i^0) &= \frac{1}{N}\sum_{i=1}^N y_i^1 - \frac{1}{N}\sum_{i=1}^Ny_i^0
%\end{align*}
%In the context of the social pressure example, this means that the average effect of the social pressure mailing is equal to the average outcome ($\frac{1}{N}\sum_{i=1}^N y_i^1$) that we would have observed if all individuals were sent the mailing, minus the average outcome ($\frac{1}{N}\sum_{i=1}^N y_i^0$) that we would have observed if all individuals were not sent the mailing (for the social pressure example, average outcomes are just turnout proportions).

%It is important to note that this decomposition is not possible for other summaries of the distribution of effects. For example if we wanted to know the median effect ($median(y_i^1- y_i^0)$), it cannot be broken up in this manner.
%\begin{align*}
%median(y_i^1- y_i^0) &\ne median(y_i^1) - median(y_i^0)
%\end{align*}
%Because the average effect can be decomposed in the above manner, and other summary statistics cannot, for the rest of this course, we will focus on average effects.\footnote{In fact, this lack of decomposability for summaries of the distribution of effects other than the average (e.g., the median effect, the lower and upper quartile of effects, the variance of effects, etc.) has important implications for the limits of causal knowledge.} In order to simplify notation, we will write our averages using expectation notation. That is, we will write $E[y^1]=\frac{1}{N}\sum_{i=1}^N y_i^1$ and $E[y^0]=\frac{1}{N}\sum_{i=1}^N y_i^0$. Therefore, the ATE can be written as $E[y^1] - E[y^0]$.

%If we denote the treatment variable (whether an individual received a mailing) with $T$, then the observed outcome can be written in terms of the potential outcomes and the treatment variable $Y_i = y_i^1 \cdot T_i + y_i^0 \cdot(1-T_i)$. Furthermore, randomization of treatment assignment justifies the ignorability assumption (GH pg. 183):
%\begin{align*} 
%y^0,y^1 \bot T
%\end{align*}
%In words, this means that the treatment assignment is independent of the potential outcomes. That is, knowing whether an individual received treatment (or control) does not tell us anything about their potential outcomes. This implies that the following equalities hold: 
%\begin{align}
%E[y^1] = E[Y|T=1] \label{eq:ig1} \\
%E[y^0] = E[Y|T=0] \label{eq:ig0}
%\end{align}
%In words, the first equation states that the expected (average) outcome in the population if \emph{everyone} received the treatment is equal to the expected (average) outcome in the population among those that \emph{did} receive the treatment, and the second equation states that the expected (average) outcome in the population if \emph{no one} received the treatment is equal to the expected (average) outcome in the population among those that \emph{did not} receive the treatment.

%We can now explain why randomized experiments allow the identification of average effects. For a large population, if we randomly assign half of the population to receive $T=1$, then the subpopulation with $T=1$ is just a random sample from the population, and if half of the population is large (think $\frac{N}{2} \rightarrow \infty$), we would expect the average $y^1$ from this subpopulation to equal the average $y^1$ in the overall population.\footnote{It is the consistency of the sample average that allows us to reach this conclusions.} Similarly, if we randomly assign half of the population to receive $T=0$, then the subpopulation with $T=0$ is just a random sample from the population, and if half of the population is large, we would expect the average $y^0$ in this subpopulation to equal the average $y^0$ in the overall population.\footnote{For a small population, the equalities in \eqref{eq:ig1} and \eqref{eq:ig0} will not hold exactly under the interpretation stated above. For example, suppose our population is the 50 states in the US, the outcome variable is statewide voter turnout, the treatment condition is having election day registration (EDR) available, and control is not having EDR available. In this case, even if we could randomly assign 25 of the states to treatment and 25 of the states to control, it is unlikely that the average voter turnout among the treated states \emph{for any given randomization} would equal the average voter turnout for all 50 states if all states received the treatment. The equalities in \eqref{eq:ig1} and \eqref{eq:ig0} will hold exactly across \emph{all possible randomizations}, but we would have to account for the additional variance across randomizations. This will be discussed in the next section. }

%More formally, if we randomize all units in the population to either $T=0$ or $T=1$, then we know that $T$ is independent from $y^1$ and $y^0$ (recall that independence between two random variables means that knowing the value of one does not provide information about the other). Independence between $T$ and $y^1$ means that the marginal distribution of $y^1$ ($f(y^1)$) is equal to the distribution of $y^1$ conditional on $T=t$ ($f(y^1|t)$). One consequence of this equality ($f(y^1) = f(y^1|t)$), is that Assumption \eqref{eq:ig1} holds. The explanation for Assumption \eqref{eq:ig0} is analogous.

%{\bf Summary:}  \\
%In order to define Individual Treatment Effects (ITEs), we must assume that SUTVA holds. To identify the Average Treatment Effect (ATE), we must also assume that the ignorability assumptions in \eqref{eq:ig1} and \eqref{eq:ig0} hold.

%{\bf Discussion Questions:} 
%\begin{itemize}
%\item How well does SUTVA seem to hold for Kuklinski et al. (1997), Gerber, Green, and Larimer (2008), and Fish (2002)
%\item How well does ignorability hold for Washington (2008)... 
%\end{itemize}



%\subsection{Estimation and Inference for Average Treatment Effects using Linear Regression}

%If SUTVA and the ignorability assumptions in \eqref{eq:ig1} and \eqref{eq:ig0} hold, then we can estimate ATE using linear regression. However, all of the assumptions we discussed in lectures 5, 6, and 7 still apply. We previously used the assumptions in lectures 5, 6, and 7 to make predictive (i.e., non-causal) inference. To make causal inference, we must first assume that SUTVA and ignorability hold. That allows us to interpret our predictive inferences as causal inferences. Then, we still have to go through the routine of making predictive inferences.  The following paragraphs summarize that process for a simple regression with a dichotomous treatment. 

%In order to obtain an unbiased estimate of ATE with linear regression, we must assume that $E[Y|T=t]$ can be modeled as a linear function of $t$.
%\begin{align}
%E[Y|T=t] = \beta_0 + \beta_1 \cdot t
%\end{align}
%This assumption is trivially satisfied when $T$ is binary. Note that when SUTVA, ignorability, and linearity of conditional expectation are satisfied, $\beta_1$ is the average treatment effect. 
%\begin{align*}
%ATE &= E[y^1 - y^0] \\
%&= E[y^1] - E[y^0] \\
%&= E[Y|T=1] - E[Y|T=0] \\
%&= [\beta_0 + \beta_1] - [\beta_0] \\
%&= \beta_1
%\end{align*}
%If we can obtain an unbiased estimate of $\beta_1$, we will have an unbiased estimate of ATE. Therefore, in addition to SUTVA, ignorability, and linearity of conditional expectation, if we have a representative sample from the population (or if our sample is the population), then the least squares estimator ($\widehat{\beta}_1$) provides an unbiased estimator of $\beta_1= ATE$.\footnote{If we have actually sampled the entire population (i.e., we have a census), then the population may not be very large and hence we may worry that the equalities stated in \eqref{eq:ig1} and \eqref{eq:ig0} will only hold across all possible randomizations (e.g., as in the example presented in fn. ?). However, it turns out that if the sample/population is not too small, treating the sample as if it had been randomly sampled from a very large population provides a good approximation to what we would get if we performed the correct finite-population inference (where the randomness comes from the randomized treatment assignment). For example, even if the population is the 50 states, treating those states as if they had been randomly sampled from some infinite superpopulation of states provides a good approximation of performing exact finite-population inference where we would take into account the randomness inherent in randomizing 25 of the states to receive EDR. We do not need to believe that the superpopulation of states is a real thing, it merely functions to provide an approximation.}







\end{document}


